\chapter{连续型随机变量}\label{ch:g12:contdist}

等待时间、物理测量值、比例：许多随机量
取值连续，任何单个值的概率都不为正。
概率则通过对\emph{概率密度}积分来计算。本
章研究均匀分布、指数分布与正态分布，并利用
正态分布刻画民意调查与样本的波动。

\section{概率密度}

\begin{definition}[概率密度，连续型随机变量]\label{def:g12:contdist:density}
区间 $I$ 上的\emph{概率密度}\index{概率密度}是指
连续、非负的函数 $f$ 定义于 $I$ 上，且 $\int_I f(t)\,\dd t = 1$，
（当 $I$ 为无界区间时，积分理解为
随有界区间不断扩大的积分的极限）。随机变量 $X$ 有
密度为 $f$，若对所有 $a \leq b$（均在 $I$ 中）：
\[
\P(a \leq X \leq b) = \int_a^b f(t)\,\dd t .
\]
\end{definition}

概率即是\emph{密度曲线下的面积}。特别地，
$\P(X = a) = \int_a^a f = 0$ 对每个单个值 $a$ 均成立，即只有区间才有概率，
且 $\P(a \leq X \leq b) = \P(a < X < b)$。

\begin{definition}[期望与方差]\label{def:g12:contdist:exp}
设 $X$ 的密度为 $f$，定义于 $I$ 上，则：
\[
\E(X) = \int_I t\,f(t)\,\dd t, \qquad
\V(X) = \int_I \bigl(t - \E(X)\bigr)^2 f(t)\,\dd t
= \E(X^2) - \E(X)^2 .
\]
\end{definition}

第 14--15 章的规则（期望的线性性、
独立变量方差的可加性、比内梅–切比雪夫不等式）仍然成立；
其证明将求和换成积分，在本阶段不加证明地承认。

\section{均匀分布}

\begin{definition}[均匀分布]\label{def:g12:contdist:uniform}
$X$ 服从\emph{均匀分布}\index{均匀分布}
$\mathcal U\intcc ab$，是指其密度为常数，
$f(t) = \frac{1}{b-a}$，定义在 $\intcc{a}{b}$ 上。于是对
$\intcc{c}{d} \subseteq \intcc ab$,
有 $\P(c \leq X \leq d) = \frac{d - c}{b - a}$，即概率与
长度成正比。
\end{definition}

\begin{proposition}\label{prop:g12:contdist:uniformmoments}
若 $X \sim \mathcal U\intcc ab$，则
$\E(X) = \dfrac{a+b}{2}$，且 $\V(X) = \dfrac{(b-a)^2}{12}$。
\end{proposition}

\begin{proof}
$\E(X) = \int_a^b \frac{t}{b-a}\dd t
= \frac{b^2 - a^2}{2(b-a)} = \frac{a+b}{2}$。对于方差，
$\E(X^2) = \frac{b^3 - a^3}{3(b-a)} = \frac{a^2 + ab + b^2}{3}$，于是
\[
\V(X) = \frac{a^2 + ab + b^2}{3} - \frac{(a+b)^2}{4}
= \frac{4a^2 + 4ab + 4b^2 - 3a^2 - 6ab - 3b^2}{12}
= \frac{(b - a)^2}{12}. \qedhere
\]
\end{proof}

\section{指数分布}

\begin{definition}[指数分布]\label{def:g12:contdist:expo}
设 $\lambda > 0$，$X$ 服从\emph{指数分布}\index{指数分布}，
$\mathcal E(\lambda)$，是指其
在 $\intco{0}{+\infty}$ 上的密度为
\[
f(t) = \lambda\,\eu^{-\lambda t}.
\]
于是 $\P(X \leq x) = 1 - \eu^{-\lambda x}$，且
$\P(X > x) = \eu^{-\lambda x}$（$x \geq 0$）。
\end{definition}

\begin{proposition}\label{prop:g12:contdist:expomoments}
若 $X \sim \mathcal E(\lambda)$，则：
$\E(X) = \dfrac{1}{\lambda}$，且 $\V(X) = \dfrac{1}{\lambda^2}$。
\end{proposition}

\begin{proof}
在 $\intcc{0}{A}$ 上分部积分：
\[
\int_0^A t\,\lambda\eu^{-\lambda t}\dd t
= \bigl[-t\,\eu^{-\lambda t}\bigr]_0^A + \int_0^A \eu^{-\lambda t}\dd t
= -A\eu^{-\lambda A} + \frac{1 - \eu^{-\lambda A}}{\lambda}
\xrightarrow[A\to+\infty]{} \frac1\lambda,
\]
利用 $A\eu^{-\lambda A} \to 0$（\cref{thm:g12:exp:growth}）。再作一次
分部积分得 $\E(X^2) = \frac{2}{\lambda^2}$，从而
$\V(X) = \frac{2}{\lambda^2} - \frac{1}{\lambda^2}$。
\end{proof}

\begin{theorem}[无记忆性]\label{thm:g12:contdist:memoryless}
\index{无记忆性}
若 $X \sim \mathcal E(\lambda)$，则对所有 $s, t \geq 0$：
\[
\pcond{X > s}{X > s + t} = \P(X > t) .
\]
指数分布刻画的是\emph{无老化}的寿命，
（如放射性原子核，而非灯泡）。
\end{theorem}

\begin{proof}
\[
\pcond{X > s}{X > s+t}
= \frac{\P(X > s + t)}{\P(X > s)}
= \frac{\eu^{-\lambda(s+t)}}{\eu^{-\lambda s}} = \eu^{-\lambda t}
= \P(X > t). \qedhere
\]
\end{proof}

\begin{omfigure}
\begin{tikzpicture}
\begin{axis}[omaxis,
  width=8.8cm, height=5cm,
  xmin=-0.4, xmax=4.4, ymin=0, ymax=1.15,
  xtick={1.3}, xticklabels={$t$}, ytick={1}, yticklabels={$\lambda$}]
  \addplot[name path=dens, omDef, thick, domain=0:4.2, samples=80]
    {exp(-x)};
  \path[name path=xaxis] (0,0) -- (4.2,0);
  \addplot[omThm!20] fill between[of=dens and xaxis,
    soft clip={domain=1.3:4.2}];
  \draw[black, dashed, thin] (axis cs:1.3,0) -- (axis cs:1.3,0.2725);
  \node[omThm, font=\small, anchor=west] at (axis cs:2.1,0.35)
    {$\P(X > t) = \eu^{-\lambda t}$};
\end{axis}
\end{tikzpicture}

{\small 指数密度 $\lambda\eu^{-\lambda x}$（此处
$\lambda = 1$）：超过 $t$ 的尾部面积（红色）为 $\eu^{-\lambda t}$，
而无记忆性表明：每段尾部都像整个分布
经重新缩放后的样子。}
\end{omfigure}

\section{正态分布}

\begin{definition}[正态分布]\label{def:g12:contdist:normal}
设 $X$ 服从\emph{标准正态分布}\index{标准正态
分布} $\mathcal N(0, 1)$，若其在 $\R$ 上的密度为
\[
\varphi(t) = \frac{1}{\sqrt{2\pi}}\,\eu^{-t^2/2}
\]
（钟形曲线）。一般地，
$X \sim \mathcal N(\mu, \sigma^2)$，若 $\dfrac{X - \mu}{\sigma} \sim
\mathcal N(0,1)$，则 $\E(X) = \mu$，$\V(X) = \sigma^2$。
\end{definition}

\begin{omfigure}
\begin{tikzpicture}
\begin{axis}[
  width=11cm, height=4.6cm,
  axis lines=middle, xmin=-4, xmax=4, ymin=0, ymax=0.45,
  xtick={-3,-2,-1,0,1,2,3},
  xticklabels={$\mu{-}3\sigma$,$\mu{-}2\sigma$,$\mu{-}\sigma$,$\mu$,
               $\mu{+}\sigma$,$\mu{+}2\sigma$,$\mu{+}3\sigma$},
  ytick=\empty, samples=120, domain=-4:4]
  \addplot[name path=gauss, omDef, thick] {exp(-x^2/2)/sqrt(2*pi)};
  \path[name path=xaxis] (-4,0) -- (4,0);
  \addplot[omDef!20] fill between[of=gauss and xaxis,
           soft clip={domain=-2:2}];
  \node[omDef] at (axis cs:0,0.17) {$\approx 95.4\%$};
\end{axis}
\end{tikzpicture}

{\small 正态密度：约 $68.3\%$ 的概率质量位于
距均值 $\sigma$ 之内，$95.4\%$ 位于 $2\sigma$ 之内，$99.7\%$ 位于
$3\sigma$ 之内。}
\end{omfigure}

\begin{remark}
因子 $\frac{1}{\sqrt{2\pi}}$ 使总面积为 $1$ --- 这是一个著名的
计算（\emph{高斯积分}），在大学中完成。
$\int \varphi$ 没有初等的原函数表达式：正态概率需
查表或借助计算器求得。
\end{remark}

\begin{theorem}[棣莫弗–拉普拉斯定理，在本阶段不加证明地承认]\label{thm:g12:contdist:tcl}
\index{棣莫弗–拉普拉斯定理}
设 $X_n \sim \mathcal B(n, p)$，且
$Z_n = \dfrac{X_n - np}{\sqrt{np(1-p)}}$（标准化二项变量）。则
对所有 $a \leq b$，
\[
\P(a \leq Z_n \leq b) \xrightarrow[n\to+\infty]{}
\int_a^b \varphi(t)\,\dd t .
\]
\emph{本结论在本阶段不加证明地承认。}它解释了钟形曲线
为何普遍出现：当 $n$ 很大时，二项分布近似
服从正态分布 --- 而（中心极限定理，属于大学内容）指出，大量微小的
独立效应之和同样近似服从正态分布。
\end{theorem}

\begin{method}[$95\%$ 水平下的波动区间]\label{met:g12:contdist:fluctuation}
\index{置信区间}
由于正态变量落在其均值 $1.96$ 个标准差之内的
概率为 $0.95$，故当 $n$ 很大时，
频率 $F_n = \frac{X_n}{n}$（$\mathcal B(n,p)$ 样本）满足
\[
\P\left(p - 1.96\sqrt{\tfrac{p(1-p)}{n}} \leq F_n \leq
p + 1.96\sqrt{\tfrac{p(1-p)}{n}}\right) \approx 0.95 .
\]
\begin{itemize}
  \item \emph{波动区间}（检验）：若假定的 $p$ 使观测频率
        落在该区间之外，则在
        $5\%$ 水平下拒绝该假设。
  \item \emph{置信区间}（估计）：简化的区间
        $\left[F_n - \frac{1}{\sqrt n},\ F_n + \frac{1}{\sqrt n}\right]$
        包含 $p$ 的概率至少为 $0.95$，
        （利用 $1.96\sqrt{p(1-p)} \leq 1$，见\cref{exo:g12:sums:8}）。
\end{itemize}
\end{method}

\begin{example}\label{ex:g12:contdist:poll}
一项对 $n = 1000$ 名选民的民意调查显示，某候选人的得票率为 $52\%$。
其置信区间 $52\% \pm \frac{1}{\sqrt{1000}} \approx 52\% \pm 3.2\%$ 仍
包含 $50\%$，因此仅凭这一民调，尚不能断定该候选人
处于领先地位。
\end{example}

\section{练习}

\begin{exercise}[$\star$]\label{exo:g12:contdist:1}
一辆公交车每隔 $15$ 分钟经过一次站点；一名旅客在
均匀随机的时刻到达。设 $X \sim \mathcal U\intcc{0}{15}$ 为等待时间，
求 $\P(X \leq 5)$，$\P(4 \leq X \leq 10)$，$\E(X)$ 及 $\sigma(X)$。
\end{exercise}

\begin{exercise}[$\star$]\label{exo:g12:contdist:2}
验证 $f(t) = \frac{3}{8}t^2$ 是 $\intcc{0}{2}$ 上的概率密度，
并求 $\E(X)$ 及 $\P(X \geq 1)$，其中 $X$ 以该函数为概率密度。
\end{exercise}

\begin{exercise}[$\star$]\label{exo:g12:contdist:3}
某电子元件的寿命（单位：年）服从
$\mathcal E(0.2)$。
\begin{enumerate}
  \item 求寿命的期望及 $\P(X > 5)$。
  \item 该元件已工作了 $3$ 年，
        求它至少还能再工作 $5$ 年的概率？
\end{enumerate}
\end{exercise}

\begin{exercise}[$\star\star$]\label{exo:g12:contdist:4}
某放射性原子的寿命服从
$\mathcal E(\lambda)$ 分布，其半衰期即中位数 $m$：$\P(X > m) = \frac12$。
将 $m$ 表示为 $\lambda$ 的函数，并与期望比较。哪一个更大？
这对偏态分布而言为何合理？
\end{exercise}

\begin{exercise}[$\star\star$]\label{exo:g12:contdist:5}
某群体的身高（单位：cm）服从 $\mathcal N(175,\ 7^2)$。利用
$68$--$95$--$99.7$ 法则，估计该群体中身高
在 $168$ 至 $182$ cm 之间、高于 $189$ cm 及低于 $154$ cm 的人所占比例。
\end{exercise}

\begin{exercise}[$\star\star$]\label{exo:g12:contdist:6}
一台机器灌装标注为 $500$ g 的袋装产品，灌装质量服从
$\mathcal N(\mu,\ 4^2)$。规定至多 $2.5\%$ 的袋子质量低于
$500$ g。利用 $95\%$ 法则，符合规定的
$\mu$ 的最小值为多少？
\end{exercise}

\begin{exercise}[$\star\star$]\label{exo:g12:contdist:7}
一颗声称均匀的骰子被掷 $1200$ 次，
其中六点出现 $260$ 次（见\cref{exo:g12:sums:7}）。
\begin{enumerate}
  \item 求六点频率的 $95\%$
        波动区间（设骰子均匀，共掷 $1200$ 次）。
  \item 观测频率是否落在区间内？比较此结论与
        比奈梅–切比雪夫分析所得结论的强弱。
\end{enumerate}
\end{exercise}

\begin{exercise}[$\star\star\star$]\label{exo:g12:contdist:8}
选举前进行民调，共访问 $n$ 人，拟估计某候选人的得票率 $p$，
估计方法为观测频率 $F_n$。
\begin{enumerate}
  \item 利用
        \cref{met:g12:contdist:fluctuation}中的置信区间，样本量至少应取多少，才能保证
        误差在 $\pm 2$ 个百分点之内？
  \item 实际中两位候选人的得票率仅相差 $1$ 个百分点。试说明
        无论调查设计得多么严谨，任何现实中的民调都无法可靠地判定胜者，
        其原因何在。
\end{enumerate}
\end{exercise}

\begin{exercise}[$\star\star\star$]\label{exo:g12:contdist:9}
设 $X$ 有概率密度 $f$，其定义区间为 $\intcc ab$，并令
$Y = \alpha X + \beta$，其中 $\alpha > 0$。
\begin{enumerate}
  \item 证明 $\P(c \leq Y \leq d)
        = \P\left(\frac{c - \beta}{\alpha} \leq X \leq
        \frac{d-\beta}{\alpha}\right)$，并由此推出 $Y$ 的概率密度为
        $g(y) = \frac{1}{\alpha} f\left(\frac{y - \beta}{\alpha}\right)$。
  \item 由此推出：若 $Z \sim \mathcal N(0,1)$，则
        $\mu + \sigma Z$ 的概率密度为
        $\frac{1}{\sigma\sqrt{2\pi}}\,
        \eu^{-(y - \mu)^2/(2\sigma^2)}$ --- 即一般正态分布的概率密度。
\end{enumerate}
\end{exercise}

\section{问题：等公交车，测人群}

\begin{problem}[{周末问题 --- 三种概率密度主宰世界：
完全无知、无记忆等待，以及众多
微小成因之钟形曲线；另有使公交车总是迟到的悖论}]
\label{pb:g12:contdist:1}
为什么你乘坐的公交车似乎总比时刻表
承诺的时间更久？为什么普通学生所在的班级
比普通班级的平均人数更多？为什么你的朋友
比你拥有更多的朋友？一个巧妙的数学思想 --- 长度偏倚
抽样 --- 可统一解答这三个问题，它出现在本章，
介于均匀分布、指数分布与正态分布之间
（\cref{def:g12:contdist:uniform}，\cref{def:g12:contdist:expo}，
\cref{def:g12:contdist:normal}）。本卷最后一题
运用这三种概率密度，并落脚于
全书的主线：钟形曲线。

\medskip
\noindent\textbf{第一部分 --- 三个角色。}
\begin{enumerate}
  \item 一辆公交车在接下来的
        $30$ 分钟内的某个均匀随机时刻到达。求 $\P(\text{等待} \leq 10)$、
        期望等待时间，以及 $\sigma$
        （\cref{prop:g12:contdist:uniformmoments}）。
  \item 验证 $f(x) = 2x$ 在 $\intcc{0}{1}$ 上是概率密度，
        并求 $\E(X)$ 及 $\P\left(X > \frac12\right)$。
  \item 热线电话的来电构成无记忆流，
        下一次来电的等待时间服从平均数为
        $10$ 分钟的指数分布。求 $\P(T > 15)$ 及中位数等待时间
        --- 为什么中位数\emph{小于}平均数？
  \item 已经等待了 $20$ 分钟，求
        $\P(T > 35 \mid T > 20)$
        （见\cref{thm:g12:contdist:memoryless}）？用一句话
        解释其含义。
  \item 将每种等待与正确的模型配对 --- 均匀分布、
        指数分布，或两者皆非：(a) 盖革计数器的咔嗒声；
        (b) 每 $8$ 分钟一班的地铁，
        你的到达时刻与班次不同步；(c) 热线电话的下一通
        来电；(d) 一只灯泡的失效，它
        \emph{会损耗}。用无记忆性定理
        论证 (d)。
\end{enumerate}

\medskip
\noindent\textbf{第二部分 --- 钟形曲线大显身手。}
（利用正态分布的标志值：$68\,\%$ 落在 $\sigma$ 之内，$95\,\%$
落在 $2\sigma$ 之内，$99.7\,\%$ 落在 $3\sigma$ 之内。）
\begin{enumerate}[resume]
  \item 成年人身高近似服从
        $\mathcal N(172,\ 8)$（厘米），求落在
        $\intcc{164}{180}$、$\intcc{156}{188}$、
        $\intcc{148}{196}$ 内的比例各是多少？
  \item 估计身高超过 $190$ 厘米的比例（先求标准分数，
        再查表或用计算器）。
  \item 智商分数经校准后服从 $\mathcal N(100, 15)$。
        求分数高于 $130$ 的比例 --- 并估计大约每
        多少人中有一人？
  \item 一台机器切割螺栓，其长度服从
        $\mathcal N(50,\ 0.1)$ 毫米，规格为
        $\intcc{49.8}{50.2}$。被拒收的比例是多少？工业界
        推崇“六西格玛”工艺，其规格限
        位于 $\pm 6\sigma$ 处：这能带来什么好处，
        为什么工厂要竭力追求它？
  \item 棣莫弗--拉普拉斯（\cref{thm:g12:contdist:tcl}）：
        对 $100$ 次公平掷硬币，近似计算
        $\P(45 \leq \text{正面} \leq 55)$（连续性修正：
        $\pm 0.5$；$\sigma = 5$）。在
        \cref{pb:g11:binom:1}中，哪台木制机器会被本定理
        平滑成一条曲线？
  \item 由\cref{exo:g12:contdist:8}可知，误差为 $\pm 2$ 个点时，
        所需样本量约为 $n = 2\,400$。求
        误差为 $\pm 0.5$ 个点所需的样本量 --- 并用一句话说明，为何
        相差一个点的竞选，
        任何民调都无法诚实地判定胜者。
\end{enumerate}

\medskip
\noindent\textbf{第三部分 --- 公交车悖论。}
\begin{enumerate}[resume]
  \item 公交车以无记忆的到达流方式到站，平均间隔为 $10$
        分钟。你在任意时刻到达车站。
        直觉认为，平均等待时间为 $5$ 分钟
        （即间隔的一半）。无记忆性定理则
        给出怎样的结论？
  \item 用一个简单的时刻表检验这一直觉：间隔交替为
        $5$ 分钟和 $15$ 分钟（平均间隔 $10$）。
        求一次均匀随机到达落在长间隔中的概率，
        再求真正的平均等待时间 ---
        并指出罪魁祸首：到达时刻选取间隔时
        概率与其\emph{长度}成正比。
  \item 在无记忆流中，你所落入的间隔的期望长度为
        $20$ 分钟 --- 是典型间隔的两倍
        （回溯至上一班车与前瞻至下一班车的时间，二者
        \emph{都}服从均值为 $10$ 的指数分布）。
        请将此结论与第 12 题协调，并用一句话叙述
        \emph{检验悖论}。
  \item 陆地上也有同样的“刺痛”：某大学开设九个班，每班
        $10$ 名学生，另有一个班 $110$ 名学生。先算平均班级规模，
        再算平均每名学生所处班级的规模，
        即\emph{平均学生}所经历的规模。宣传册会引用哪个数字，
        哪个才是学生们真实的体验？
  \item 友谊悖论 --- “你的朋友平均而言比你拥有更多的
        朋友”，其背后是同一数学。
        用一句话说明，什么东西扮演了
        “长间隔”的角色。
\end{enumerate}

\medskip
\noindent\textbf{第四部分 --- 本册结语。}
\begin{enumerate}[resume]
  \item 模拟：从业者的桥梁。若 $U$ 在
        $\intoo{0}{1}$ 上服从均匀分布，证明
        $T = -10\ln U$ 服从均值为 $10$ 的指数分布（计算
        $\P(T > t)$）。业界所有随机等待模拟
        都依赖这一行小技巧。
  \item 一个古老的游戏程序员配方可构造出钟形曲线：将
        十二个独立的 $\intcc{0}{1}$ 上均匀分布的随机变量相加，
        再减去 $6$。求该和的平均数与方差，
        并结合高尔顿板 --- 解释为什么
        结果近似服从正态分布。
  \item 市场崩盘比钟形曲线所允许的更为剧烈：1987 年的
        崩盘被称为“$25\sigma$ 事件”，若按正态假设，
        其概率约为 $10^{-137}$。
        正确的结论是什么 --- 是关于世界，还是关于
        模型？（一句话，统计学家的
        信条。）
  \item 结语，并向本册告别：三个角色各用一行概括，
        均匀分布（有界范围内的无知）、
        指数分布（无记忆的危险率）、
        正态分布（许多微小原因之和）；检验
        悖论作为本章的“刺痛”；以及整套丛书的
        脉络 --- 从六年级数纸箱，到
        测量每一群人的钟形曲线 --- 而大学分册
        正等待着极限、
        积分与中心极限定理接手之处。
\end{enumerate}
\end{problem}
